\section{总结与延伸}

整场讲座可以压缩为一个双层结论。关于模型：它们会组合抽象 feature，在一次 forward pass 中并行运行多条计算路径，并让未来目标提前约束当前 token；这些能力解释了低资源迁移、医学组合、多语言共享、加法、押韵等现象，也解释了自述不忠实、幻觉与 jailbreak 中的内部冲突。关于方法：稀疏替换、attribution graph 与 original-model intervention 能提供局部因果证据，但 reconstruction error、feature ambiguity、success-case bias 与 unexplained attention 决定了它仍是一台不完美的科学仪器。

\begin{importantbox}{一句话总纲}
不要问“模型是否真的在思考”这种不可操作的问题；要问在给定 prompt 上，哪些内部表示被激活、哪些路径并行贡献、哪些未来目标提前出现，以及怎样用干预证伪这套解释。
\end{importantbox}

\subsection{课堂结论与实践清单}

\small
\begin{multicols}{2}
\begin{enumerate}
  \item 行为正确不等于机制单一或稳定；
  \item 单 neuron 通常不是最佳语义单位；
  \item Feature label 必须与因果作用一起解释；
  \item CLT 提升可读性，也可能损失保真度；
  \item Reconstruction error 应显式进入图与报告；
  \item Attention 未解释时，不要声称已解释策略选择；
  \item 抽象表示可以跨语言和任务复用；
  \item 一次 forward pass 内可有大规模并行计算；
  \item 模型自述不保证忠实反映内部算法；
  \item IDK、回答和 refusal 可能是竞争路径；
  \item Planning 需要“提前 feature + 干预改变轨迹”证据；
  \item 安全监控应覆盖生成过程，而非只看输入；
  \item 机制结论必须报告 prompt、模型与时间范围；
  \item 失败案例与成功案例同样属于方法结果。
\end{enumerate}
\end{multicols}
\normalsize

\enlargethispage{4\baselineskip}
\subsection{拓展阅读}

\small
\noindent\textbf{互动案例：}\href{https://transformer-circuits.pub/2025/attribution-graphs/biology.html}{On the Biology of a Large Language Model}，包含多步推理、诗歌规划、多语言、加法、医学诊断、幻觉、拒答与 CoT faithfulness 的完整交互图。

\noindent\textbf{方法论文：}\href{https://transformer-circuits.pub/2025/attribution-graphs/methods.html}{Circuit Tracing: Revealing Computational Graphs in Language Models}，详细讨论 replacement model、Cross-Layer Transcoder、graph pruning、评估与局限。\textbf{阅读顺序：}先用本讲义建立 feature--graph--intervention 的证据阶梯，再回到互动论文逐项点击节点；遇到漂亮图时，优先寻找 reconstruction error、干预结果、失败案例和 attention 限定。
\normalsize

\end{document}
